\documentclass[10pt,a4paper]{amsart}
\usepackage{amsmath,amssymb,amsthm,graphicx}
% Retained minimal preamble for the Stacks Project chapter more-algebra.tex (More on Algebra).
% The chapter file pinned-original-source.tex carries only \input{preamble} on its first line, so no
% preamble text lives inside the chapter itself.  The Stacks Project uses one single shared preamble for
% every chapter; the snapshot stages it identically as sources/stacks-chow/preamble.tex and
% sources/stacks-etale-cohomology/preamble.tex.  The block below is copied mechanically, line for line,
% from sources/stacks-chow/preamble.tex lines 14-18 (the Xy-pic setup) and 154-217 (the native theorem
% declarations together with the mathematical macros).  Nothing was added, removed or redefined; the
% wrapper class and the amsmath/amssymb/amsthm/graphicx base packages are supplied by the assembler.
% Chapter: more-algebra.tex -- More on Algebra.
% For commutative diagrams we use Xy-pic
\usepackage[all]{xy}

% We use 2cell for 2-commutative diagrams.
\xyoption{2cell}

% Theorem environments.
%
\theoremstyle{plain}
\newtheorem{theorem}[subsection]{Theorem}
\newtheorem{proposition}[subsection]{Proposition}
\newtheorem{lemma}[subsection]{Lemma}

\theoremstyle{definition}
\newtheorem{definition}[subsection]{Definition}
\newtheorem{example}[subsection]{Example}
\newtheorem{exercise}[subsection]{Exercise}
\newtheorem{situation}[subsection]{Situation}

\theoremstyle{remark}
\newtheorem{remark}[subsection]{Remark}
\newtheorem{remarks}[subsection]{Remarks}

\numberwithin{equation}{subsection}

% Macros
%
\def\lim{\mathop{\mathrm{lim}}\nolimits}
\def\colim{\mathop{\mathrm{colim}}\nolimits}
\def\Spec{\mathop{\mathrm{Spec}}}
\def\Hom{\mathop{\mathrm{Hom}}\nolimits}
\def\Ext{\mathop{\mathrm{Ext}}\nolimits}
\def\SheafHom{\mathop{\mathcal{H}\!\mathit{om}}\nolimits}
\def\SheafExt{\mathop{\mathcal{E}\!\mathit{xt}}\nolimits}
\def\Sch{\mathit{Sch}}
\def\Mor{\mathop{\mathrm{Mor}}\nolimits}
\def\Ob{\mathop{\mathrm{Ob}}\nolimits}
\def\Sh{\mathop{\mathit{Sh}}\nolimits}
\def\NL{\mathop{N\!L}\nolimits}
\def\CH{\mathop{\mathrm{CH}}\nolimits}
\def\proetale{{pro\text{-}\acute{e}tale}}
\def\etale{{\acute{e}tale}}
\def\QCoh{\mathit{QCoh}}
\def\Ker{\mathop{\mathrm{Ker}}}
\def\Im{\mathop{\mathrm{Im}}}
\def\Coker{\mathop{\mathrm{Coker}}}
\def\Coim{\mathop{\mathrm{Coim}}}

% Boxtimes
%
\DeclareMathSymbol{\boxtimes}{\mathbin}{AMSa}{"02}

%
% Macros for moduli stacks/spaces
%
\def\QCohstack{\mathcal{QC}\!\mathit{oh}}
\def\Cohstack{\mathcal{C}\!\mathit{oh}}
\def\Spacesstack{\mathcal{S}\!\mathit{paces}}
\def\Quotfunctor{\mathrm{Quot}}
\def\Hilbfunctor{\mathrm{Hilb}}
\def\Curvesstack{\mathcal{C}\!\mathit{urves}}
\def\Polarizedstack{\mathcal{P}\!\mathit{olarized}}
\def\Complexesstack{\mathcal{C}\!\mathit{omplexes}}
% \Pic is the operator that assigns to X its picard group, usage \Pic(X)
% \Picardstack_{X/B} denotes the Picard stack of X over B
% \Picardfunctor_{X/B} denotes the Picard functor of X over B
\def\Pic{\mathop{\mathrm{Pic}}\nolimits}
\def\Picardstack{\mathcal{P}\!\mathit{ic}}
\def\Picardfunctor{\mathrm{Pic}}
\def\Deformationcategory{\mathcal{D}\!\mathit{ef}}

% --- appended 2026-10-12 (Chief): the project own definitions of the reference/slogan/history
% display environments, copied mechanically from the shared Stacks preamble lines 8-12
% (\usepackage{verbatim} plus the three \newenvironment lines, which the chapter preamble keeps
% verbatim from the project preamble). Nothing added, removed or redefined: these environments
% carry no mathematical content and the project itself typesets them as comments. Without these
% three lines a counted statement whose own text uses them fails with Environment slogan undefined.
\usepackage{verbatim}
\newenvironment{reference}{\comment}{\endcomment}
\newenvironment{slogan}{\comment}{\endcomment}
\newenvironment{history}{\comment}{\endcomment}
\usepackage{cleveref}
\crefname{originallemma}{original lemma}{original lemmas}
\Crefname{originallemma}{Original lemma}{Original lemmas}
\makeatletter
\newlabel{lemma-pull-pseudo-coherent}{{0650}{source}{}{sourcepointer-1}{}}
\newlabel{lemma-pull-pseudo-coherent@cref}{{[originallemma][0][]0650}{[1][0][]source}{}{}{}}
\makeatother
\crefname{originallemma}{original lemma}{original lemmas}
\Crefname{originallemma}{Original lemma}{Original lemmas}
\makeatletter
\newlabel{lemma-lift-complex-stably-frees}{{0BCB}{source}{}{sourcepointer-2}{}}
\newlabel{lemma-lift-complex-stably-frees@cref}{{[originallemma][0][]0BCB}{[1][0][]source}{}{}{}}
\makeatother
\crefname{originallemma}{original lemma}{original lemmas}
\Crefname{originallemma}{Original lemma}{Original lemmas}
\makeatletter
\newlabel{derived-lemma-morphisms-from-projective-complex}{{064B}{source}{}{sourcepointer-3}{}}
\newlabel{derived-lemma-morphisms-from-projective-complex@cref}{{[originallemma][0][]064B}{[1][0][]source}{}{}{}}
\makeatother
\title{A pseudo-coherent complex over a local ring is determined up to isomorphism by the dimensions of its residue field fibres}
\author{The Stacks Project Authors}
\date{stacks-project:more-algebra.tex:local-snapshot-20261012; GNU FDL 1.2 or later}
\begin{document}
\maketitle
\setcounter{subsection}{1}
\noindent\textit{Editorial note.} All delivered mathematical text is the author's own, copied from the pinned Stacks Project chapter file more-algebra.tex; no statement or proof step was rewritten, reordered or removed. The counted statement is the labeled lemma environment at source lines 19547-19563 (canonical tag 0BCC) and the counted proof is the proof environment at source lines 19565-19591, which immediately follows it and is its own proof.\par\medskip
\noindent\textit{Editorial note.} The statement is written in two sentences both beginning with Then, as printed by the author; the text is copied exactly as it stands and is not corrected.\par\medskip
\noindent\textit{Editorial note.} The preamble span runs to line 90, the current end of the retained preamble, which includes the appended environment definitions; the counted text does not use them, and the span is taken whole so that the file version is unambiguous.\par\medskip
\noindent\textit{Editorial note.} The source URL is the source-level Stacks Project URL, as required when more than one card is produced from one source; the per-claim tag is carried by claim\_key and by this note. The counted statement and proof cite no bibliography entry, so this card needs no bibliography.\par\medskip
\noindent\textit{Allowed standard prerequisites (editorial).}\begin{itemize}
\item Pseudo-coherence is preserved by pulling back along the residue field map, so the fibre of the complex at the residue field is pseudo-coherent and its cohomology spaces are finite dimensional and vanish above some degree. invoked at proof lines 19565-19569. Establishes that the dimensions in the statement are finite and eventually zero, which is what the displayed complex uses.
\item Every object of the derived category of a field is isomorphic to a complex with zero differentials, whose terms are the cohomology spaces. used at proof lines 19569-19572. Puts the fibre in the normal form to which the lifting statement applies.
\item A map in the derived category between bounded complexes of projectives is represented by a map of complexes, and a map of finite free modules whose reduction modulo the maximal ideal is an isomorphism is an isomorphism. used at proof lines 19574-19591. Supplies the map of complexes in the uniqueness argument and upgrades its degreewise reductions to isomorphisms.
\end{itemize}\medskip
\ifdefined\hypertarget\hypertarget{sourcepointer-1}{}\fi
\noindent\textit{Source reference.} Original lemma 0650: sources/stacks-more-algebra/pinned-original-source.tex, source lines 16345--16345. The reference is applied at proof lines 19565-19566 to see that the fibre at the residue field is pseudo-coherent, which gives finiteness and vanishing of its cohomology; the rest of the existence argument is inside the counted span.\par\medskip
\ifdefined\hypertarget\hypertarget{sourcepointer-2}{}\fi
\noindent\textit{Source reference.} Original lemma 0BCB: sources/stacks-more-algebra/pinned-original-source.tex, source lines 19517--19517. The reference produces the complex of free modules representing the object from the fibre data, applied at proof lines 19572-19573; the uniqueness half of the proof is inside the counted span.\par\medskip
\ifdefined\hypertarget\hypertarget{sourcepointer-3}{}\fi
\noindent\textit{Source reference.} Original lemma 064B: sources/stacks-chow/tags.txt, source lines 7860--7860. The reference lets the isomorphism between the two representatives be realized by a map of complexes, invoked at proof lines 19579-19582; the degreewise reduction argument that follows is inside the counted span.\par\medskip
\begin{lemma}
\label{lemma-lift-pseudo-coherent-from-residue-field}
Let $(R, \mathfrak m, \kappa)$ be a local ring. Let $K \in D(R)$
be pseudo-coherent. Set
$d_i = \dim_\kappa H^i(K \otimes_R^\mathbf{L} \kappa)$.
Then $d_i < \infty$ and for some $b \in \mathbf{Z}$ we have
$d_i = 0$ for $i > b$.
Then there exists a complex
$$
\ldots \to
R^{\oplus d_{b - 2}} \to
R^{\oplus d_{b - 1}} \to
R^{\oplus d_b} \to 0 \to \ldots
$$
representing $K$ in $D(R)$. Moreover, this complex is unique up to
isomorphism(!).
\end{lemma}

\begin{proof}
Observe that $K \otimes_R^\mathbf{L} \kappa$ is pseudo-coherent
as an object of $D(\kappa)$, see Lemma \ref{lemma-pull-pseudo-coherent}.
Hence the cohomology spaces are finite dimensional and vanish above
some cutoff. Every object of $D(\kappa)$ is isomorphic
in $D(\kappa)$ to a complex $E^\bullet$ with zero differentials.
In particular $E^i \cong \kappa^{\oplus d_i}$ is finite free.
Applying Lemma \ref{lemma-lift-complex-stably-frees} we obtain
the existence.

\medskip\noindent
If we have two complexes $F^\bullet$ and $G^\bullet$
with $F^i$ and $G^i$ free of rank $d_i$ representing $K$.
Then we may choose a map of complexes $\beta : F^\bullet \to G^\bullet$
representing the isomorphism $F^\bullet \cong K \cong G^\bullet$, see
Derived Categories, Lemma
\ref{derived-lemma-morphisms-from-projective-complex}.
The induced map of complexes
$\beta \otimes 1 : F^\bullet \otimes_R^\mathbf{L} \kappa \to
G^\bullet \otimes_R^\mathbf{L} \kappa$
must be an isomorphism of complexes as the differentials
in $F^\bullet \otimes_R^\mathbf{L} \kappa$ and
$G^\bullet \otimes_R^\mathbf{L} \kappa$ are zero.
Thus $\beta^i : F^i \to G^i$ is a map of finite free $R$-modules
whose reduction modulo $\mathfrak m$ is an isomorphism.
Hence $\beta^i$ is an isomorphism and we win.
\end{proof}

\end{document}
